% !TeX root = ./../main.tex
% Słownik pojęć.
% Hasło pojawia się w słowniku dopiero po użyciu w tekście: \gls{etykieta}, \glspl{etykieta}.

\newglossaryentry{qubit}{
  name={kubit},
  plural={kubity},
  sort={kubit},
  description={
    Podstawowa jednostka informacji kwantowej.
    Jej stan jest opisywany znormalizowanym wektorem $\ket{\psi}=\alpha\ket{0}+\beta\ket{1}$ w dwuwymiarowej zespolonej przestrzeni Hilberta, gdzie $\alpha,\beta\in\mathbb{C}$ oraz $|\alpha|^2+|\beta|^2=1$.
  }
}

\newglossaryentry{decoherence}{
  name={dekoherencja},
  sort={dekoherencja},
  description={
    Utrata kwantowych własności układu wskutek niekontrolowanego oddziaływania z otoczeniem.
  }
}

\newglossaryentry{nisq}{
  name={era NISQ},
  sort={era NISQ},
  description={
    Okres rozwoju komputerów kwantowych o średniej skali i znacznym poziomie szumu (ang. \emph{Noisy Intermediate-Scale Quantum}), w którym pełna korekcja błędów nie jest jeszcze dostępna.
  }
}

\newglossaryentry{error-correction}{
  name={korekcja błędów kwantowych},
  sort={korekcja bledow kwantowych},
  description={
    Grupa metod chroniących informację kwantową przez zakodowanie kubitów logicznych w nadmiarowej liczbie kubitów fizycznych oraz wykrywanie błędów pomiarem syndromu i odwracanie ich w trakcie obliczeń (ang. \emph{quantum error correction}, QEC).
  }
}

\newglossaryentry{error-mitigation}{
  name={mitygacja błędów kwantowych},
  sort={mitygacja bledow kwantowych},
  description={
    Grupa metod zmniejszających wywołane szumem obciążenie wartości oczekiwanych lub rozkładów wyników przez klasyczną obróbkę wyników wielu wykonań obwodu, jego zmodyfikowanych wariantów lub obwodów kalibracyjnych, bez kodowania informacji w kubitach logicznych (ang. \emph{quantum error mitigation}, QEM).
  }
}

\newglossaryentry{error-reduction}{
  name={redukcja błędów kwantowych},
  sort={redukcja bledow kwantowych},
  description={
    W niniejszej pracy: grupa metod korygujących wynik obliczeń kwantowych wyłącznie po stronie klasycznej, modelem korekty wyznaczonym eksperymentalnie dla danego zadania i urządzenia, bez dodatkowych kubitów i bramek w obwodzie zadania; dodatkowe wykonania obwodu służą tylko kalibracji modelu.
  }
}

\newglossaryentry{logical-qubit}{
  name={kubit logiczny},
  plural={kubity logiczne},
  sort={kubit logiczny},
  description={
    Kubit zakodowany w stanie wielu kubitów fizycznych za pomocą kodu korekcyjnego; pełni w obliczeniach rolę pojedynczego kubitu chronionego przed błędami.
  }
}

\newglossaryentry{syndrome}{
  name={syndrom błędu},
  sort={syndrom bledu},
  description={
    Wynik pomiaru wykonywanego w kodzie korekcyjnym, który wskazuje rodzaj i położenie błędu, lecz nie ujawnia zakodowanej informacji.
  }
}

\newglossaryentry{quantum-channel}{
  name={kanał kwantowy},
  sort={kanal kwantowy},
  description={
    Liniowe odwzorowanie macierzy gęstości, które jest całkowicie dodatnie i zachowuje ślad (CPTP); najogólniejszy opis zmiany stanu układu kwantowego, obejmujący bramki unitarne i szum.
  }
}

\newglossaryentry{coherent-error}{
  name={błąd koherentny},
  sort={blad koherentny},
  description={
    Błąd będący stałą, niezamierzoną operacją unitarną, na przykład obrotem o dodatkowy kąt wskutek niedokładnej kalibracji; jest jednakowy w każdym wykonaniu obwodu i nie zamienia stanu czystego w mieszany.
  }
}

\newglossaryentry{incoherent-error}{
  name={błąd niekoherentny},
  sort={blad niekoherentny},
  description={
    Błąd stochastyczny: kanał, którego nie da się zapisać jedną operacją unitarną, na przykład losowe działanie bramek Pauliego, relaksacja lub defazowanie; zamienia stany czyste w mieszane.
  }
}

\newglossaryentry{unital-channel}{
  name={kanał unitalny},
  sort={kanal unitalny},
  description={
    Kanał kwantowy, który nie zmienia stanu maksymalnie mieszanego; dla jednego kubitu nie przesuwa środka kuli Blocha. Kanał nieunitalny, na przykład relaksacja, ściąga układ do wyróżnionego stanu.
  }
}

\newglossaryentry{crosstalk}{
  name={przesłuch},
  sort={przesluch},
  description={
    Niezamierzone oddziaływanie między kubitami lub liniami sterującymi (ang. \emph{crosstalk}), wskutek którego operacja na jednym kubicie zaburza inne kubity albo błąd bramki zależy od operacji wykonywanych równocześnie w sąsiedztwie.
  }
}

\newglossaryentry{qubit-leakage}{
  name={wyciek poza przestrzeń obliczeniową},
  sort={wyciek poza przestrzen obliczeniowa},
  description={
    Przejście układu fizycznego realizującego kubit do stanu spoza podprzestrzeni rozpiętej przez $\ket{0}$ i $\ket{1}$ (ang. \emph{leakage}), na przykład na wyższy poziom energetyczny.
  }
}

\newglossaryentry{systematic-error}{
  name={składowa systematyczna błędu},
  sort={skladowa systematyczna bledu},
  description={
    Różnica między prawdopodobieństwem wyniku na zaszumionym urządzeniu a prawdopodobieństwem idealnym; nie zależy od liczby pomiarów i powtarza się w kolejnych seriach, dopóki parametry urządzenia są ustalone.
  }
}

\newglossaryentry{shot-noise}{
  name={szum próbkowania},
  sort={szum probkowania},
  description={
    Składowa statystyczna błędu wyniku (ang. \emph{shot noise}, także: szum śrutowy): losowa różnica między częstością otrzymaną ze skończonej liczby pomiarów a prawdopodobieństwem wyniku; jej odchylenie standardowe maleje jak odwrotność pierwiastka z liczby pomiarów.
  }
}

\newglossaryentry{distribution-shrinkage}{
  name={współczynnik skurczenia rozkładu},
  sort={wspolczynnik skurczenia rozkladu},
  description={
    Liczba $s\in[0,1]$, przez którą szum niekoherentny mnoży odchylenia rozkładu wyników od rozkładu jednostajnego; dla kanału depolaryzującego działającego na cały rejestr $s=1-\lambda$.
  }
}

\newglossaryentry{ptm}{
  name={macierz transferu Pauliego},
  sort={macierz transferu Pauliego},
  description={
    Rzeczywista macierz kanału kwantowego w bazie operatorów Pauliego (ang. \emph{Pauli Transfer Matrix}, PTM); dla jednego kubitu składa się z pierwszego wiersza (zachowanie śladu), wektora przesunięcia (część nieunitalna) i bloku unitalnego $3\times3$ (obrót i skrócenie wektora Blocha).
  }
}

\newglossaryentry{quantum-sampling}{
  name={próbkowanie kwantowe},
  sort={probkowanie kwantowe},
  description={
    Sposób prowadzenia obliczeń kwantowych, w którym wynikiem jest rozkład prawdopodobieństwa stanów bazowych wyznaczony z wielokrotnego wykonania tego samego obwodu.
  }
}

\newglossaryentry{qcosamp}{
  name={QCoSamp},
  sort={QCoSamp},
  description={
    Kwantowe próbkowanie szeregiem cosinusów (ang. \emph{Cosine Series Quantum Sampling}): metoda budowy operatorów, dla których rozkład wyników pomiaru jednego kubitu jest ważoną sumą składników $\frac{1}{2}(1+\cos(nx+r))$ argumentu $x$ wprowadzonego do obwodu; pochodzi z pracy Wereszczyńskiego i~in.
  }
}

\newglossaryentry{pdo}{
  name={operator zniekształcenia fazy},
  sort={operator znieksztalcenia fazy},
  description={
    (ang. \emph{Phase Distortion Operator}, PDO) Operator unitarny diagonalny w bazie obliczeniowej, który mnoży amplitudę każdego stanu bazowego przez czynnik fazowy zależny od kodowanych danych; służy do zapisu danych w fazach lokalnych.
  }
}

\newglossaryentry{pdu-function}{
  name={funkcja PDU},
  sort={funkcja PDU},
  description={
    Funkcja odwracania zniekształceń fazy (ang. \emph{Phase Distortion Unraveling}): interpolacja błędów wyznaczonych eksperymentalnie jako różnice między zmierzonymi a oczekiwanymi wartościami funkcji referencyjnej; odejmuje się ją od wyników późniejszych obliczeń wykonywanych na tym samym urządzeniu i układzie kubitów.
  }
}

\newglossaryentry{reference-function}{
  name={funkcja referencyjna},
  sort={funkcja referencyjna},
  description={
    Funkcja realizowana przez obwód kwantowy w modelu próbkowania, której wartości w węzłach siatki kalibracyjnej są znane niezależnie od tego obwodu; w niniejszej pracy składnik cosinusowy $\frac{1}{2}(1+\cos x)$ metody QCoSamp.
  }
}

\newglossaryentry{qimr}{
  name={kwantowa reprezentacja obrazu},
  plural={kwantowe reprezentacje obrazów},
  sort={kwantowa reprezentacja obrazu},
  description={
    Sposób zapisania obrazu cyfrowego w stanie rejestru kwantowego wraz z procedurą odtworzenia intensywności pikseli z wyników pomiarów, np. FRQI, NEQR lub LPIQE.
  }
}

\newglossaryentry{superposition}{
  name={superpozycja},
  sort={superpozycja},
  description={
    Stan kwantowy będący kombinacją liniową co najmniej dwóch stanów bazowych o niezerowych amplitudach zespolonych; amplitudy mogą się wzajemnie wzmacniać lub znosić.
  }
}

\newglossaryentry{entanglement}{
  name={splątanie},
  sort={splatanie},
  description={
    Własność stanu układu złożonego, którego nie da się zapisać jako iloczyn tensorowy stanów podukładów.
  }
}

\newglossaryentry{global-phase}{
  name={faza globalna},
  sort={faza globalna},
  description={
    Czynnik $\mathrm{e}^{\mathrm{i}\chi}$ wspólny dla wszystkich składowych wektora stanu; nie wpływa na wyniki żadnego pomiaru.
  }
}

\newglossaryentry{local-phase}{
  name={faza lokalna},
  sort={faza lokalna},
  description={
    Faza względna składowych superpozycji, np. $\varphi$ w stanie $(\ket{0}+\mathrm{e}^{\mathrm{i}\varphi}\ket{1})/\sqrt{2}$; niewidoczna w pomiarze w bazie obliczeniowej, staje się mierzalna dzięki interferencji.
  }
}

\newglossaryentry{interference}{
  name={interferencja kwantowa},
  sort={interferencja kwantowa},
  description={
    Sumowanie się amplitud różnych ścieżek obliczenia prowadzących do tego samego wyniku; prawdopodobieństwo wyniku zależy od różnicy faz tych amplitud.
  }
}

\newglossaryentry{bloch-sphere}{
  name={sfera Blocha},
  sort={sfera Blocha},
  description={
    Geometryczna reprezentacja stanu kubitu jako punktu sfery jednostkowej (stany czyste) lub kuli (stany mieszane); faza lokalna jest azymutem tego punktu.
  }
}

\newglossaryentry{density-matrix}{
  name={macierz gęstości},
  sort={macierz gestosci},
  description={
    Hermitowska, dodatnio półokreślona macierz o śladzie 1 opisująca stan kwantowy, także mieszany; jej diagonala zawiera prawdopodobieństwa stanów bazowych, a elementy pozadiagonalne -- koherencje.
  }
}

\newglossaryentry{transpilation}{
  name={transpilacja},
  sort={transpilacja},
  description={
    Przekształcenie obwodu logicznego w równoważny obwód złożony z bramek natywnych urządzenia i zgodny z topologią połączeń jego kubitów; zmienia liczbę i rodzaj bramek, a przez to błąd wyniku.
  }
}

\newglossaryentry{phase-kickback}{
  name={przeniesienie fazy},
  sort={przeniesienie fazy},
  description={
    (ang. \emph{phase kickback}) Efekt działania bramki kontrolowanej na kubit docelowy w stanie własnym jej operatora: wartość własna pojawia się jako faza lokalna kubitu kontrolnego.
  }
}

\newglossaryentry{uncomputation}{
  name={dekomputacja},
  sort={dekomputacja},
  description={
    (ang. \emph{uncomputation}) Odwrócenie części obliczenia po skopiowaniu jego wyniku, przywracające kubity pomocnicze do stanu początkowego, tak aby wyniki pośrednie nie blokowały interferencji.
  }
}

\newglossaryentry{overfitting}{
  name={przeuczenie},
  sort={przeuczenie},
  description={
    (ang. \emph{overfitting}) Dopasowanie modelu do szczegółów i szumu zbioru uczącego zamiast do zależności ogólnych, widoczne jako malejąca strata na zbiorze uczącym przy rosnącej stracie na zbiorze walidacyjnym.
  }
}

\newglossaryentry{data-leakage}{
  name={przeciek danych},
  sort={przeciek danych},
  description={
    (ang. \emph{data leakage}) Wpływ informacji ze zbioru, na którym ocenia się model, na jego uczenie lub wybór, prowadzący do zawyżonej oceny jakości modelu.
  }
}

\newglossaryentry{hybrid-model}{
  name={model hybrydowy kwantowo-klasyczny},
  sort={model hybrydowy kwantowo-klasyczny},
  description={
    Model uczenia maszynowego, w którym obwód kwantowy i klasyczne sieci lub funkcje tworzą jeden graf obliczeń uczony wspólnie; w niniejszej pracy sieć klasyczna wyznacza kąty bramek obwodu o stałej strukturze.
  }
}

\newglossaryentry{learned-encoding}{
  name={kodowanie uczone},
  sort={kodowanie uczone},
  description={
    Kwantowa reprezentacja obrazu, w której stanu nie zadaje wzór, lecz sieć klasyczna nauczona na zbiorze obrazów: sieć wyznacza kąty bramek obwodu o stałej strukturze, a intensywności pikseli odczytuje się z prawdopodobieństw stanów bazowych; kodowanie jest przybliżone także bez szumu.
  }
}

\newglossaryentry{hellinger-fidelity}{
  name={wierność Hellingera},
  sort={wiernosc Hellingera},
  description={
    Miara zgodności dwóch rozkładów prawdopodobieństwa wyników pomiaru, $\bigl(\sum_x\sqrt{p(x)\,q(x)}\bigr)^2\in[0,1]$, równa 1 tylko dla rozkładów identycznych; jest wiernością stanów kwantowych diagonalnych w bazie pomiaru.
  }
}

\newglossaryentry{equivariance}{
  name={ekwiwariancja},
  sort={ekwiwariancja},
  description={
    Własność odwzorowania, które przekształcenie wejścia przez element grupy przenosi na znane przekształcenie wyjścia: $f(T_g\mathbf{x})=T'_g f(\mathbf{x})$ dla każdego elementu $g$ grupy.
  }
}
